% =====================================================================
% Bagian 5 — Validasi Pengguna dan Kelayakan Adopsi  (target 1,0 halaman)
% Detail didorong ke Lampiran D (heuristik) dan E (biaya).
% Subbagian Validasi Ahli: \input{template_validasi_ahli} setelah 5.1.
% =====================================================================

\section{Validasi Pengguna dan Kelayakan Adopsi}

\subsection{Uji pengguna dan iterasi yang dihasilkannya}

Sebelum uji dengan pengguna, antarmuka melewati 19 putaran evaluasi heuristik
\cite{nielsen1990} dan audit aksesibilitas terinstrumentasi yang menghasilkan
sepuluh perubahan; empat di antaranya menyangkut cacat yang hanya terdeteksi
lewat pengukuran, kegagalan ambang kontras WCAG dan pemisahan warna di bawah
$\Delta E = 4{,}3$ pada protanopia (Lampiran~\ref{app:heuristik}). Evaluasi
heuristik adalah metode yang sah, tetapi ia \textbf{melengkapi, bukan
menggantikan} uji dengan pengguna, dan kami tidak menyajikannya sebagai
demikian.

Uji pengguna dijalankan pada tiga partisipan di luar tim yang belum pernah
melihat produk ini, dengan metode \emph{think-aloud} dan lima tugas tanpa
petunjuk cara mengerjakan. Dua tugas, menentukan \emph{alert} prioritas dan
menandai keputusan, lancar pada ketiganya sehingga tidak diubah. Tiga tugas
lain memunculkan kesulitan, dan ketiganya menghasilkan perubahan antarmuka yang
tercatat pada log putaran kerja (Tabel~\ref{tab:iterasi}).

\begin{table}[!htbp]
\caption{Temuan uji pengguna dan perubahan antarmuka yang dihasilkannya.}
\label{tab:iterasi}
\centering
\small
\begin{tabular}{p{2.9cm}p{4.2cm}p{4.6cm}}
\hline
\textbf{Kesulitan} & \textbf{Bukti} & \textbf{Perubahan} \\
\hline
Membedakan sumber pemicu di peta &
Satu partisipan gagal total tanpa bantuan setelah 35 detik dan empat salah
klik; satu lainnya berhasil namun lambat &
Legenda warna tetap ditambahkan di peta \\
Memahami keadaan data yang belum dianalisis &
Satu partisipan menyangka aplikasi rusak dan berpindah halaman &
Tanda "---" diganti label eksplisit beserta ikon, warna netral yang berbeda
dari warna galat \\
Menafsirkan lencana modalitas &
Dua partisipan menjawab benar namun ragu; satu sempat keliru lalu mengoreksi
diri &
Tooltip singkat ditambahkan pada lencana modalitas \\
\hline
\end{tabular}
\end{table}

Kegagalan total diprioritaskan di atas kesulitan yang sekadar memperlambat.
Keterbatasannya kami catat pada Bagian~\ref{sec:keterbatasan}: tiga partisipan
bukan operator posko, dan pengujian dilakukan tanpa tekanan waktu.

% ---------------------------------------------------------------------
% SISIPKAN DI SINI:  \input{template_validasi_ahli}
% ---------------------------------------------------------------------

\subsection{Kelayakan adopsi dan struktur biaya}

Sistem dirancang menambah lapisan konfirmasi pada alur kerja yang sudah berjalan,
bukan menggantikannya, dan dioperasikan oleh operator posko, bukan tenaga khusus
pembelajaran mesin. Jalur adopsi empat tahap, aspek kelayakan implementasi
beserta status verifikasinya, dan tiga dasar skalabilitas yang terukur diuraikan
pada Lampiran~\ref{app:adopsi}. Ketahanan modalitas berkontribusi pada kelayakan adopsi,
tetapi perlu dirumuskan dengan tepat: narasumber ahli menyatakan drone berkamera
inframerah sudah banyak dipakai, sehingga argumen "armada eksisting mayoritas
tanpa sensor termal" \textbf{tidak kami pakai}. Rumusan yang didukung bukti
adalah bahwa ketahanan modalitas menurunkan syarat perangkat keras minimum dan
menjaga sistem tetap memberi keluaran ketika sensor termal mati, gagal, atau
tidak terpasang pada sebagian armada.

Tiga sifat struktur biaya lebih penting daripada angka absolutnya
(Lampiran~\ref{app:biaya}): tidak ada biaya per-inferensi karena model berjalan
di perangkat sendiri, sehingga menambah frekuensi patroli tidak menambah biaya
perangkat lunak sama sekali; biaya berulang satu-satunya adalah hosting, dan itu
pun nol pada skala pilot; dan belanja modal, perangkat tepi serta sensor
termal, bersifat opsional, bukan syarat untuk mulai.

Seluruh fitur yang diklaim dapat diperiksa langsung pada \emph{deployment} yang
berjalan: klasifikasi, skor keandalan modalitas, dan peta lokalisasi berasal dari
keluaran model nyata yang ditandai otomatis oleh penanda sumber data, sementara
lapisan hotspot menarik data satelit langsung, pada satu verifikasi, 2.467 titik
terdeteksi dan 400 titik ber-FRP tertinggi disajikan dengan jumlah total tetap
dinyatakan (Lampiran~\ref{app:firms}).
